% method.tex
% Method section. Describes the SF-GRU architecture this paper evaluates
% (originally proposed by Rasouli et al., BMVC 2019 -- see
% rasouli2020sfgru in refs.bib), the JAAD labeling artifact this paper
% identifies and corrects, and the domain-adversarial + contrastive
% training protocol that forms this paper's core positive result.

\section{Method}
\label{sec:method}

\subsection{Base Architecture: SF-GRU}
\label{sec:method-sfgru}

We evaluate SF-GRU (Stacked-Fusion GRU)~\cite{rasouli2020sfgru}, a
pedestrian crossing-intention architecture originally proposed by
Rasouli~et~al. SF-GRU processes each input modality with its own
single-layer GRU, concatenates that GRU's output sequence with the raw
features of the next modality, and feeds the result into the next GRU in
the stack; the final GRU's last hidden state is passed through a linear
layer and sigmoid to produce a crossing probability. We use the
following modality order, matching the original paper and its released
implementation: local cropped-pedestrian appearance
(\texttt{local\_box}), local surrounding scene context
(\texttt{local\_context}), 2D body pose (\texttt{pose}), and raw
bounding-box trajectory (\texttt{box}). \texttt{local\_box} and
\texttt{local\_context} are 512-dimensional per-frame features pooled
from a frozen, ImageNet-pretrained VGG16 backbone; \texttt{pose} is a
36-dimensional per-frame keypoint vector (18 keypoints $\times$ 2D
coordinates) from RTMPose~\cite{rtmpose}; \texttt{box} is a 4-dimensional
bounding-box coordinate delta. We exclude ego-vehicle speed from our
main results (Section~\ref{sec:method-crossprotocol}) and separately
test a synthesized speed signal for JAAD
(Section~\ref{sec:method-speed}). All modalities are observed over a
15-frame window sampled 60 frames before the annotated crossing/non-crossing
event, matching the original protocol's observation and
time-to-event parameters.

We use a PyTorch reimplementation of the original Keras/TensorFlow SF-GRU
that preserves the model architecture, data-preparation logic (frame
sampling, class-balancing via horizontal-flip augmentation of the
minority class on the training split only, bounding-box normalization),
and training protocol (AdamW, focal binary cross-entropy loss, an
$F_1$-maximizing decision threshold selected on a held-out validation
split), differing from the original only in framework and in using
RTMPose in place of OpenPose for keypoint extraction.

\subsection{Datasets}
\label{sec:method-data}

We use the two datasets SF-GRU was originally evaluated
on: PIE~\cite{rasouli2019pie} (6 hours of continuous ego-vehicle
footage, Toronto, Canada; standard split: train~=~set01+02+04,
val~=~set05+06, test~=~set03) and
JAAD~\cite{rasouli2017jaad} (346 short clips, 2 vehicles, 3 dashcam
models, 5 countries; default split with an official train/val/test
partition). Both datasets provide per-frame bounding boxes and a
per-pedestrian crossing-intention label; we follow each dataset's
standard \texttt{seq\_type=`crossing'} sequence-generation protocol.

\subsection{A Labeling Artifact in JAAD's Default Configuration}
\label{sec:method-labelbug}

JAAD's pedestrian identifiers fall into three categories, distinguishable
by a substring in the pedestrian ID: bystanders (identifiers containing
\texttt{`p'}), pedestrians with full behavioral annotation
(identifiers containing \texttt{`b'}), and a remaining category with
neither marker. The dataset loader's \texttt{\_get\_crossing()} function
excludes bystanders explicitly, but for pedestrians in the third,
unmarked category it hard-codes both the crossing label and the
observation endpoint (\texttt{acts\,=\,[[0]]},
\texttt{event\_frame\,=\,-1}) regardless of that pedestrian's actual
behavior, because no behavioral annotation exists to determine it. The
dataset loader's public interface exposes a \texttt{sample\_type}
parameter (default \texttt{`all'}) that controls whether these
unmarked-category pedestrians are included; the standard, published
convention -- matching the \emph{JAAD\_behavior} subset used to report
the 686-instance pedestrian count cited in prior
work~\cite{kotseruba2021pcpa} -- is \texttt{sample\_type=`beh'},
restricting to only the fully behaviorally-annotated (\texttt{`b'})
category.

Because \texttt{sample\_type} defaults to \texttt{`all'} and no
downstream training script in the codebase we built on overrides it, every
JAAD-side experiment we initially ran silently trained and evaluated
against a label distribution in which 72--80\% of samples per split
(depending on split) carry a label that is not a behavioral observation
but a fallback constant. This is not merely severe class imbalance: the
label for the majority of samples is definitionally uninformative,
independent of any true crossing behavior. We discovered this by
noticing an anomalously low, near-constant positive rate (9.2--9.9\%
across splits) that did not match the literature's reported
JAAD\_behavior base rates, then confirmed it by direct inspection of the
dataset loader's source and cross-checked with an independent
literature audit. Correcting to \texttt{sample\_type=`beh'} raises the
positive rate to 62--83\% across splits (train 82.6\%, val 73.9\%, test
62.6\%) and reduces the usable JAAD training set from several hundred
samples to 195 -- a real, unavoidable cost of the fix, since the
excluded pedestrians genuinely lack behavioral ground truth and cannot
be relabeled. PIE has no equivalent unmarked-pedestrian category and is
unaffected. \emph{All results in this paper use the corrected
\texttt{sample\_type=`beh'} configuration.} We flag this as a
methodological finding independent of our other results: because this
default is silent and the resulting label distribution is not
obviously degenerate (a near-10\% positive rate is plausible on its
face for a crossing-intention task), any prior work built on this
dataset loader without explicitly setting \texttt{sample\_type} is at
risk of the same corruption, and we recommend the community audit
downstream JAAD-based pipelines for this default.

\subsection{Cross-Dataset Evaluation Protocol}
\label{sec:method-crossprotocol}

The original SF-GRU evaluation, like nearly all pedestrian
crossing-intention work we are aware of prior to
Gesnouin~et~al.~\cite{gesnouin2022} and including all 2024--2026 work
surveyed in Section~\ref{sec:related}, trains and tests within a single
dataset. To assess generalization under dataset shift, we additionally
adopt the cross-dataset protocol of~\cite{gesnouin2022}: a model is
trained (and its decision threshold tuned) on one dataset's train/val
split, then evaluated, unchanged, on the \emph{other} dataset's held-out
test split. We report both directions, PIE~$\to$~JAAD and
JAAD~$\to$~PIE.

Because PIE's \texttt{speed} channel is a continuous ego-vehicle OBD
sensor reading with no JAAD equivalent -- the original implementation
substitutes a discrete \texttt{vehicle\_act} encoding in JAAD's
\texttt{speed} slot -- we exclude \texttt{speed} entirely for our main
cross-dataset experiments, following the same exclusion
in~\cite{gesnouin2022}, to avoid conflating a feature-semantics mismatch
with genuine domain shift. Our main results therefore use models
trained on \texttt{\{local\_box, local\_context, pose, box\}}.

\subsection{Ruling Out Camera Geometry}
\label{sec:method-geometry}

A natural hypothesis is that PIE and JAAD's differing camera
installations (PIE: single vehicle, fixed, published intrinsic and
extrinsic calibration, including fisheye distortion coefficients;
JAAD: multiple vehicles, no published calibration) explain cross-dataset
failure through a recoverable geometric mismatch rather than a learned
one. We test this directly by projecting each pedestrian's foot point to
ground-plane metric coordinates using PIE's real published calibration
(camera height 1.270\,m, pitch $-10^\circ$, full intrinsic matrix and
fisheye distortion model) and researched estimates for JAAD (height
1.40\,m, pitch $-5^\circ$, intrinsics rescaled from PIE's to JAAD's
frame width), then retraining SF-GRU on the resulting metric-space
trajectory instead of raw pixel deltas. We validate the projection
independently by checking that the implied real-world pedestrian speed
distribution is physically plausible (mean 0.32\,m/s, 100\% within
$[0,3]$\,m/s) before using it for training. We additionally sweep
JAAD's unknown camera height ($1.2$--$1.6$\,m) and pitch ($-3^\circ$ to
$-10^\circ$) over 20 grid points, reusing a single trained checkpoint's
cached features and varying only the ground-plane projection at
evaluation time, to test whether \emph{any} plausible camera-parameter
choice for JAAD changes the outcome (Section~\ref{sec:results-geometry}).

\subsection{Domain-Adversarial Training}
\label{sec:method-dann}

Motivated by ruling out geometry (Section~\ref{sec:results-geometry})
and a per-modality zeroing ablation
(Section~\ref{sec:results-ablation}) that localizes cross-dataset
sensitivity to the visual context branch rather than the raw box
trajectory, we train SF-GRU under an unsupervised domain-adaptation
(UDA) protocol: a Domain-Adversarial Neural Network
(DANN)~\cite{ganin2015dann}. A domain classifier, reached from the
model's final fused hidden state through a Gradient Reversal Layer
(GRL), is trained to predict which dataset (source or target) a sample
came from; the main network is trained, via the reversed gradient, to
make that prediction as difficult as possible. The source dataset is
fully supervised (crossing labels and domain label); the target
dataset's train-split \emph{inputs only} are used for the domain
adversary, its labels and held-out val/test splits never touched during
training -- the same train-on-source, evaluate-unchanged-on-target
protocol as our other cross-dataset experiments, with the addition of
unlabeled target inputs available at training time for the domain
classifier.

We additionally combine the GRL with a contrastive scale-invariance
term: for each source-domain training batch, we synthesize a second
``view'' of each sample by randomly rescaling its box trajectory around
a fixed center (simulating the same physical motion filmed at a
different distance or focal length), extract the box modality's own GRU
hidden state for both views, and apply an NT-Xent contrastive loss
pulling the two views' embeddings together while pushing apart other
samples in the batch. This trains the box embedding to encode motion
pattern rather than absolute pixel scale, independently of the GRL's
broader domain-adversarial pressure on the fused representation. The
total training loss is $\mathcal{L} = \mathcal{L}_{\text{crossing}} +
\lambda_{\text{dom}} \mathcal{L}_{\text{domain}} +
\lambda_{\text{con}} \mathcal{L}_{\text{contrastive}}$, with
$\lambda_{\text{dom}}=1.0$ and the DANN gradient-reversal coefficient
following the standard sigmoid ramp schedule of~\cite{ganin2015dann}. We
report both the a~priori default contrastive hyperparameters
($\lambda_{\text{con}}=0.5$, NT-Xent temperature $0.2$) and a small
follow-up grid search over $\lambda_{\text{con}} \in
\{0.25,0.5,1.0\}$, temperature $\in \{0.1,0.2,0.4\}$
(Section~\ref{sec:results-dann}), confirming any selected configuration
at $n=6$ seeds before reporting it as a headline number.

\subsection{Context-Targeted Domain Adversarial Training}
\label{sec:method-context}

The recipe above applies the GRL to the final fused representation,
after all four modalities have been concatenated through the stack.
Motivated by the modality-zeroing ablation (Section~\ref{sec:results-ablation}),
which localizes cross-dataset sensitivity to the \texttt{local\_context}
modality specifically, we also test a variant in which the GRL and
domain classifier instead read only the \texttt{local\_context}
branch's own hidden state, tapped immediately after its GRU and before
concatenation with the next modality -- a read-only tap that does not
alter the main classification forward path. The contrastive term is
unchanged, still operating on the box branch. This tests whether
concentrating the domain-adversarial pressure at the diagnosed leakage
source matches, exceeds, or falls short of regularizing the whole fused
representation; same protocol otherwise (default contrastive
hyperparameters, $n=6$ seeds per direction,
Section~\ref{sec:results-context}).

\subsection{Representational Invariance Probe}
\label{sec:method-probe}

To test whether the GRL's domain-adversarial training achieves the
representational invariance it is designed to produce, we take a
trained checkpoint's final fused hidden state on held-out PIE and JAAD
test data and fit a \emph{separate}, freshly-initialized classifier --
never touching the original model's weights -- to predict which dataset
each representation came from. We report both a linear probe (logistic
regression) and a small non-linear probe (2-layer MLP, 32 hidden units,
early stopping) to rule out a linear probe being too weak to detect
invariant structure a stronger probe would find, averaged over 5
independent random train/test splits of the probe's own data. This is
deliberately distinct from the GRL's own in-training domain-accuracy
diagnostic, which is trained adversarially against a simultaneously
changing feature extractor and is not a reliable measure of the final,
frozen representation's actual separability
(Section~\ref{sec:results-probe}). We additionally report a trivial
control -- a probe given only each sample's track length, a single
scalar carrying no pedestrian-behavior information -- to establish how
separable PIE and JAAD are from a dataset-fingerprint statistic alone.

\subsection{Second-Architecture Replication}
\label{sec:method-crossattn}

All results above use SF-GRU's stacked, sequential fusion (each
modality's GRU output concatenated into the next modality's input,
Section~\ref{sec:method-sfgru}). To test whether the GRL + contrastive
recipe's effect and the probe dissociation are specific to this fusion
mechanism, we implement the same GRL + contrastive training recipe
(Section~\ref{sec:method-dann}) on a structurally distinct fusion
design: box-anchored cross-attention. Each modality is encoded
independently by its own GRU (returning the full hidden-state
sequence, not just the final state); the box modality's encoded
sequence serves as the attention query, the concatenation of every
other modality's encoded sequence serves as the key/value context, and
the attended output is passed through one further GRU before the
classification head -- mirroring the box-anchored design used in
prior cross-modal-attention variants of SF-GRU, but, to our knowledge,
never previously evaluated under the cross-dataset protocol. The GRL
and NT-Xent contrastive taps are placed at the analogous points to
the stacked-GRU model: the GRL reads the final post-attention GRU's
hidden state (the same representation the classification head reads),
and the contrastive term reads the box modality's own encoder output.
We train this architecture at the a~priori default contrastive
hyperparameters ($\lambda_{\text{con}}=0.5$, temperature $0.2$,
Section~\ref{sec:method-dann}) rather than re-sweeping for this
architecture specifically, since the goal is to test whether the
recipe's effect generalizes, not to find its separately-optimal
hyperparameters here; $n=6$ seeds per direction, matching our other
headline confirmations.

\subsection{Third-Architecture Replication}
\label{sec:method-uncertainty}

We additionally test a third fusion mechanism: uncertainty-weighted
fusion, in which each modality is encoded independently by its own GRU
into a single final hidden state (no attention or sequential
concatenation), a per-modality log-variance is predicted from that
hidden state by a small linear head, and the modalities are combined by
a softmax over negative log-variance (an inverse-variance,
confidence-weighted average) before the classification head. We choose
this architecture specifically, rather than an arbitrary third choice,
because it is the strongest \emph{unadapted} baseline in an 8-way
cross-dataset architecture sweep we ran across every fusion mechanism
implemented in this codebase (PIE~$\to$~JAAD AUC $0.585\pm0.009$,
$n=3$, versus SF-GRU's $0.498\pm0.009$) -- making it the best-motivated
test of whether domain-adversarial training adds value on top of an
architecture that already transfers comparatively well without any
adaptation. The GRL and contrastive taps read, respectively, the fused
(post-uncertainty-weighting) representation and the box modality's own
independent encoder output (available directly here, since each
modality's encoder is independent rather than embedded mid-sequence in
a shared stack). Same protocol as Section~\ref{sec:method-crossattn}:
default contrastive hyperparameters, $n=6$ seeds per direction.

\subsection{Pretrained Trajectory Representation}
\label{sec:method-pretrained-traj}

Motivated by Gesnouin~et~al.'s~\cite{gesnouin2022} finding that generic
pretrained visual backbones generalize better under PIE/JAAD domain
shift than task-specific ones, we test whether the same holds for the
box-trajectory modality specifically. We pretrain a small Transformer
encoder (2 layers, 4 attention heads, 64-dimensional hidden size) on a
masked-frame reconstruction objective: for each box-delta sequence in a
pooled, unlabeled corpus (PIE's and JAAD's train and validation splits
combined, crossing labels never touched), a random subset of frames is
replaced with a learned mask token and the encoder must reconstruct
each masked frame's true 4-dimensional delta from the unmasked context,
trained with mean-squared-error loss on masked positions only. After
pretraining, we freeze the encoder and splice its per-timestep output
into SF-GRU's box GRU slot in place of the raw box delta, training the
box GRU (now consuming a 64-dimensional input) and everything
downstream of it under the same GRL + contrastive recipe and default
hyperparameters as Section~\ref{sec:method-dann}, with the pretrained
encoder's weights held fixed throughout. We note explicitly that this
is a weaker test of the motivating hypothesis than an ideal version
would be, in two independent respects: our pretraining corpus (pooled
PIE and JAAD trajectories) is far smaller and less diverse than the
external, large-scale motion corpora (e.g.\ Waymo Open Motion,
Argoverse~2) a genuine test would use, \emph{and} it is drawn from the
same two datasets used for downstream cross-dataset evaluation, unlike
Gesnouin~et~al.'s source/target-independent visual backbones. We did
not have local access to an external motion corpus, and report this as
a locally-feasible approximation of the intended experiment, not an
equivalent of it (Section~\ref{sec:results-pretrained-traj}).

\subsection{Fine-Tuned Vision-Language Model}
\label{sec:method-vlm}

As a further test of whether pretrained world-knowledge aids
cross-dataset transfer -- this time using a backbone genuinely
pretrained on a large, source/target-independent corpus, unlike
Section~\ref{sec:method-pretrained-traj}'s locally-pretrained
trajectory encoder -- we LoRA-fine-tune~\cite{hu2022lora}
Qwen2-VL-2B-Instruct~\cite{qwen2vl} (rank 16, targeting all attention
and MLP projection matrices) on a VQA-style formulation: the model is
shown the observation frame with the pedestrian's bounding box drawn
on it and a fixed prompt asking whether the pedestrian will cross,
trained to generate ``YES''/``NO'' via next-token cross-entropy on the
answer span only. We train separately on each dataset's train split
(minority-class oversampled to balance, matching this paper's
\texttt{balanced=True} convention elsewhere) and evaluate with the same
train-on-source/test-on-target protocol as every other result in this
paper. Because the model is queried via greedy decoding, evaluation
yields a hard YES/NO prediction rather than a calibrated probability;
we report accuracy, precision, recall, and F1, and do not claim these
are directly comparable to the AUC figures reported elsewhere in this
paper.

Checkpoint selection for the JAAD-trained direction requires care:
JAAD's validation split under \texttt{sample\_type=`beh'} is only 17
sequences (the same constraint noted for JAAD-as-source training
throughout this paper), making a single ``best-by-validation-accuracy''
checkpoint choice high-variance. Rather than select one checkpoint, we
save the three best checkpoints by validation accuracy during training
and evaluate each independently on the full cross-dataset test split,
reporting the distribution across them rather than a single point
estimate (Section~\ref{sec:results-vlm}).

\subsection{Synthesizing a Speed Modality for JAAD}
\label{sec:method-speed}

PIE's \texttt{speed} channel is real, continuous ego-vehicle OBD sensor
data; JAAD has no equivalent signal, so every experiment above excludes
speed from both datasets to keep the modality set comparable, meaning
PIE's model never uses a real signal it natively has. We test whether
recovering this modality for JAAD helps, by estimating a coarse
per-frame ego-motion proxy from JAAD's raw video via dense optical flow
(Farneback), computed over the region \emph{excluding} the pedestrian's
own bounding box (background motion approximates ego-vehicle motion)
and reduced to the per-frame median flow magnitude. We z-score each
dataset's speed channel independently using training-split statistics
before use, since PIE's real km/h-scale readings and JAAD's unitless
optical-flow magnitude are not on a comparable absolute scale. This
gives both datasets the full five-modality set
(\texttt{local\_box, local\_context, pose, box, speed}) rather than the
four-modality set used elsewhere in this paper
(Section~\ref{sec:results-speed}).
